% Problem record op_38f9dbb6c3f92b0c. This TeX file is derived from
% database/problems_json/op_38f9dbb6c3f92b0c.json by scripts/sync-tex.mjs; edit the JSON record
% and rerun the script rather than editing this file.
\section{Strong converse for the quantum capacity of a general channel}
\label{sec:op_38f9dbb6c3f92b0c}

\paragraph{Problem.}
Does every finite-dimensional memoryless quantum channel obey the strong converse at its unassisted quantum capacity? Let $\mathcal N:\mathcal L(A)\to\mathcal L(B)$ be a completely positive trace-preserving map. An $n$-use entanglement-transmission code consists of arbitrary quantum channels $\mathcal E_n:\mathcal L(S_n)\to\mathcal L(A^{\otimes n})$ and $\mathcal D_n:\mathcal L(B^{\otimes n})\to\mathcal L(\widehat S_n)$, where $\dim S_n=\dim\widehat S_n=M_n$. The parties share no assisting resource and exchange no classical messages. With $\Phi_{M_n}$ the normalized maximally entangled state between a reference $R_n$ and the message system, define
\begin{equation}
 F_n:=\operatorname{Tr}\!\left[\Phi_{M_n}^{R_n\widehat S_n}(\operatorname{id}_{R_n}\otimes\mathcal D_n\circ\mathcal N^{\otimes n}\circ\mathcal E_n)(\Phi_{M_n}^{R_n S_n})\right],\qquad r_n:=\frac1n\log_2 M_n.
 \label{eq:h13-code}
\end{equation}
The quantum capacity $Q(\mathcal N)$ is the supremum of $\liminf_n r_n$ over sequences of codes in Eq.~\eqref{eq:h13-code} with $F_n\to1$. Is it true that every code sequence satisfies
\begin{equation}
 \liminf_{n\to\infty}r_n>Q(\mathcal N)\quad\Longrightarrow\quad F_n\longrightarrow0?
 \label{eq:h13-converse}
\end{equation}
A proof must cover arbitrary encoders and collective decoders. A counterexample must give one fixed finite-dimensional channel and a sequence violating Eq.~\eqref{eq:h13-converse}.

\subsection*{Status}

Solved

\subsection*{Source}

Hayashi discusses the general strong-converse gap in Section~9.9.3, pp.~554--555, after presenting the Rains-information converse bound in Theorem~9.12 \sourcecite{ref:h13-book}{Hay17}. Tomamichel, Wilde, and Winter formulate the general question in their Introduction and distinguish the quantum capacity from its strong-converse capacity \sourcecite{ref:h13-tww}{TWW17}.

\subsection*{Progress}

\begin{itemize}
  \item The regularized coherent information determines the vanishing-error threshold $Q(\mathcal N)$ but does not imply Eq.~\eqref{eq:h13-converse}; see Hayashi's Theorem~9.10 and Section~9.9.3 \sourcecite{ref:h13-book}{Hay17}.

  \item Tomamichel, Wilde, and Winter's Theorem~8 gives an exponential strong-converse upper bound using the regularized channel Rains information, even with classical pre- and post-processing. Their Section~VI proves equality with capacity for generalized dephasing channels. The general upper bound need not coincide with $Q(\mathcal N)$ \sourcecite{ref:h13-tww}{TWW17}.

  \item Kondra, Brinster, Kampermann, Bru{\ss}, and Wyderka prove an all-code exponential strong converse for finite-dimensional antidegradable channels in Theorem~1 and degradable channels in Theorem~4 of their August~2026 preprint. This includes erasure channels over the full parameter range. Their Theorem~11 also covers a specified nondegradable multilevel amplitude-damping family. Their general-channel bound in Theorem~8 does not determine the capacity, as discussed in Section~VIII \sourcecite{ref:h13-kondra}{KBK+26}.

  \item Tomamichel proves an exponential strong converse for Pauli channels within the class of stabilizer codes, with arbitrary decoders; see Section~6 of the July~2026 preprint. The restriction concerns the encoded subspace and cannot simply be removed. Section~7, Problem~7.3, identifies the then-open extension to unrestricted encoders \sourcecite{ref:h13-tomamichel}{Tom26}.

  \item Beigi and Tomamichel, Theorem 1, prove an exponential strong converse for every finite-dimensional memoryless quantum channel: at every fixed rate gap $\gamma>0$ above $Q(\mathcal N)$ there is $\alpha_\gamma>0$ such that the entanglement-transmission fidelity obeys $F_n<2^{-\alpha_\gamma n}$ for all sufficiently large $n$. The model permits mixed encoded states and arbitrary joint decoders. This proves the archived implication in Eq.~\eqref{eq:h13-converse}, in the stronger exponential form \sourcecite{ref:h13-beigi-tomamichel}{BT26}.

  \item Cheng and Tomamichel, Theorem 5, give a strictly positive entanglement-generation strong-converse exponent at every rate above $Q(\mathcal N)$. Their Appendix D transfers exponential decay to entanglement transmission; the exact generation and transmission exponents need not coincide. Their proof uses an Arimoto approach and R\'enyi continuity bounds, whereas Beigi--Tomamichel use a fully quantum blowing-up lemma \sourcecite{ref:h13-cheng-tomamichel}{CT26}.
\end{itemize}

\subsection*{References}

\begin{enumerate}[label={},leftmargin=4.8em,labelsep=0.5em]
  \footnotesize
  \item[\textup{[Hay17]}]\label{ref:h13-book}
  M. Hayashi, \emph{Quantum Information Theory: Mathematical Foundation}, 2nd ed., Graduate Texts in Physics, Springer (2017). \href{https://doi.org/10.1007/978-3-662-49725-8}{doi:10.1007/978-3-662-49725-8}.

  \item[\textup{[TWW17]}]\label{ref:h13-tww}
  M. Tomamichel, M. M. Wilde, and A. Winter, ``Strong Converse Rates for Quantum Communication,'' \emph{IEEE Transactions on Information Theory} \textbf{63}(1), 715--727 (2017). \href{https://arxiv.org/abs/1406.2946}{arXiv:1406.2946}.

  \item[\textup{[KBK+26]}]\label{ref:h13-kondra}
  T. V. Kondra, R. Brinster, H. Kampermann, D. Bru{\ss}, and N. Wyderka, ``Sharp Quantum Capacity Thresholds: Exponential Strong Converses for Degradable and Antidegradable Channels,'' preprint (2026). \href{https://arxiv.org/abs/2608.01308}{arXiv:2608.01308}.

  \item[\textup{[Tom26]}]\label{ref:h13-tomamichel}
  M. Tomamichel, ``A strong converse for stabilizer codes over Pauli channels via the blowing-up lemma,'' preprint (2026). \href{https://arxiv.org/abs/2607.23450}{arXiv:2607.23450}.

  \item[\textup{[BT26]}]\label{ref:h13-beigi-tomamichel}
  S. Beigi and M. Tomamichel, ``Strong Converse for Quantum Capacity via a Fully Quantum Blowing-Up Lemma,'' arXiv preprint, version 1, 10 September 2026. \href{https://arxiv.org/abs/2609.11771v1}{arXiv:2609.11771v1}.

  \item[\textup{[CT26]}]\label{ref:h13-cheng-tomamichel}
  H.-C. Cheng and M. Tomamichel, ``No information transmission through quantum channels above capacity,'' arXiv preprint, version 1, 8 September 2026. \href{https://arxiv.org/abs/2609.08998v1}{arXiv:2609.08998v1}.
\end{enumerate}

\subsection*{Comment}

Updated on 2026-09-14. The complete resolution rests on two September 2026 version-1 arXiv preprints, neither identified as peer-reviewed at this audit. Their methods differ, but both papers have Tomamichel as an author, so they are not independent confirmations by disjoint author teams. Both disclose assistance with their mathematical arguments. The Solved status follows the catalog policy of recording complete preprint resolutions with an explicit publication caveat; it does not certify an independent refereeing of either technical proof. This record asks for convergence to zero, without requiring an exponential rate. The general exponential theorem also settles the transpose-degradable subclass, whose capacity is single-letter. The separate question of whether strictly transpose-degradable channels exist is unchanged. The degradable-channel record retains the earlier subclass resolution. Rains-information achievability under an enlarged PPT assistance model is a separate operational question. In particular, the book's Rains-relative-entropy bound must not be identified with max-Rains information merely because it is named an SDP bound.

\subsection*{Field}

Quantum Communication

\subsection*{Topic}

Quantum capacity; Strong converses

\subsection*{ID}

\texttt{op\_38f9dbb6c3f92b0c}
